\begin{abstract}
Single-cell RNA sequencing requires representations that separate cell
populations despite sparse expression and heterogeneous local neighborhoods.
Masked autoencoders exploit within-cell gene context, while leaving local
cell--cell structure unused. We introduce scVICAR, which recovers each original
cell from a bounded neighborhood-mixed view while retaining the clean scMAE
objective. Two implementations separate the shared principle from topology
adaptation: scVICAR-F uses fixed mixing, whereas scVICAR-T combines
topology-informed edge weights with a cell-specific mixing coefficient. The
operator preserves a local convex hull, bounds input and representation
perturbation, and corresponds to one step of nonuniform graph diffusion.
Its contamination bound links displacement to cross-class neighbor mass and
motivates controlled edge-replacement tests. Across a 15-dataset primary
comparison of 16 method configurations, scVICAR-T and scVICAR-F ranked first and second on
all four mean metrics, improving ARI over scMAE by 5.08\% and 3.93\%. Under a
separately frozen six-dataset protocol, scVICAR-T achieved the highest mean ARI
among the compared methods (0.617), exceeding Original scMAE by 2.13\%.
Matched-backbone effects were smaller and heterogeneous, indicating that the
effect of local mixing depends on the dataset. Controlled edge replacement showed that topology-informed affinity
identified injected cross-class edges (mean AUROC 0.918), while the cell-specific
coefficient tracked neighborhood purity (mean Spearman 0.520). Low-label probes
improved across most datasets, with relative gains of 0.13--0.34\%; marker-based
endpoints varied by dataset. These results support graph-vicinal anchor recovery
as a competitive local regularizer and show that topology-adaptive effects on
clean-graph clustering remain dataset dependent.
\end{abstract}
